% !TeX root = main.tex

\chapter{Schr\"odinger Bridges and Non-Conservative Transport}

\section{Motivation}

The previous chapter introduced entropy, Gibbs measures, and variational
principles. We now use those ideas to move from deterministic optimal
transport to stochastic transport.

The central object of this chapter is the \emph{Schr\"odinger bridge}. At a
high level, a Schr\"odinger bridge asks for the most likely stochastic
evolution connecting two prescribed probability distributions. In modern
generative modeling language, it asks for a probability path

\[
\rho_0=\rho_a,
\qquad
\rho_1=\rho_b,
\]

together with a velocity or drift field that transports samples from the
initial marginal to the terminal marginal.

The new feature is that the path is not chosen only by shortest transport
cost. It is regularized by diffusion, entropy, or a reference stochastic
process. This makes Schr\"odinger bridges a natural meeting point of:

\begin{itemize}
    \item optimal transport,
    \item stochastic control,
    \item diffusion models,
    \item multi-marginal interpolation,
    \item policy improvement under entropy regularization.
\end{itemize}

For our long-term goal, the important question is slightly sharper:

\[
\text{What happens if the bridge is not energy-conserving?}
\]

This question is the doorway to contact geometry.

\section{From Optimal Transport to Schr\"odinger Bridges}

\subsection{The deterministic transport picture}

In the Benamou--Brenier formulation, the squared Wasserstein distance is
obtained by minimizing kinetic energy:

\[
W_2^2(\rho_a,\rho_b)
=
\inf_{\rho_t,v_t}
\int_0^1
\int_{\R^d}
\|v_t(x)\|^2
\rho_t(x)\,dx\,dt,
\]

subject to

\[
\partial_t\rho_t+\nabla\cdot(\rho_t v_t)=0,
\qquad
\rho_0=\rho_a,
\qquad
\rho_1=\rho_b.
\]

The continuity equation says that probability mass is transported by the
velocity field $v_t$. There is no creation or destruction of mass.

\subsection{Adding diffusion}

A Schr\"odinger bridge replaces the deterministic continuity equation by a
diffusive equation of Fokker--Planck type:

\[
\partial_t\rho_t+\nabla\cdot(\rho_t v_t)
=
\varepsilon \Delta \rho_t.
\]

The diffusion coefficient $\varepsilon>0$ measures the strength of the
stochastic perturbation. The velocity $v_t$ now controls a noisy transport
process rather than a purely deterministic one.

\begin{remark}
When $\varepsilon\to 0$, the Schr\"odinger bridge problem formally approaches
optimal transport. For positive $\varepsilon$, the bridge prefers stochastic
paths that are close to a reference diffusion.
\end{remark}

\section{The Generalized Schr\"odinger Bridge}

The version that will matter for us includes a potential energy $U$. The
action is

\[
J(\rho,v)
=
\int_0^1
\int_{\R^d}
\left(
\frac12\|v_t(x)\|^2+U(x)
\right)
\rho_t(x)\,dx\,dt.
\]

The first term is kinetic energy. The second term penalizes or rewards regions
of state space through the potential $U$.

\begin{definition}[Generalized Schr\"odinger bridge]
A generalized Schr\"odinger bridge between $\rho_a$ and $\rho_b$ is a solution
of

\[
\inf_{\rho_t,v_t}
\int_0^1
\int_{\R^d}
\left(
\frac12\|v_t(x)\|^2+U(x)
\right)
\rho_t(x)\,dx\,dt
\]

subject to

\[
\partial_t\rho_t+\nabla\cdot(\rho_t v_t)
=
\varepsilon\Delta\rho_t,
\qquad
\rho_0=\rho_a,
\qquad
\rho_1=\rho_b.
\]
\end{definition}

One may also impose intermediate marginal constraints

\[
\rho_{t_m}=\rho_m,
\qquad
m=1,\dots,M.
\]

This is important for applications where one observes snapshots of a system at
several times but not the full underlying trajectory.

\section{Hamiltonian Structure}

The optimality conditions of the generalized bridge can be written in
Hamiltonian form. Introduce a scalar field $S_t(x)$, which plays the role of a
momentum or cotangent variable. At the optimum, the velocity is a gradient:

\[
v_t(x)=\nabla S_t(x).
\]

The density and momentum then evolve according to a coupled system of the form

\[
\partial_t\rho_t
=
-\nabla\cdot(\rho_t\nabla S_t),
\]

and

\[
\partial_t S_t
=
-\frac12\|\nabla S_t\|^2
+\frac12\varepsilon^2
\frac{\delta I}{\delta\rho}(\rho_t)
+U.
\]

Here

\[
I(\rho)
=
\int_{\R^d}
\|\nabla\log\rho(x)\|^2
\rho(x)\,dx
\]

is the Fisher information.

\begin{remark}
The field $S_t$ should be read as the cotangent representative of the motion.
The density derivative $\partial_t\rho_t$ is tangent to the probability
manifold, while $S_t$ is the dual object whose gradient gives the velocity
field moving the density.
\end{remark}

The corresponding Hamiltonian is naturally split into kinetic and potential
parts:

\[
H(\rho,S)
=
\underbrace{
\frac12
\int
\|\nabla S(x)\|^2\rho(x)\,dx
}_{K(\rho,S)}
+
\underbrace{
F(\rho)
}_{\text{potential contribution}}.
\]

The Fisher information belongs to the potential contribution. It is the trace
left by stochastic diffusion after rewriting the bridge in Hamiltonian
variables.

\section{The Jacobi Metric Viewpoint}

Classical mechanics has a useful geometric reformulation: under suitable
conditions, Hamiltonian trajectories can be seen as geodesics for a rescaled
metric called a Jacobi metric.

In the Wasserstein setting, the base metric is the Wasserstein metric
$g_{W_2}$. The generalized bridge can be viewed as a geodesic problem for a
conformal rescaling

\[
g_J
=
\bigl(H-F(\rho)\bigr)g_{W_2}.
\]

The factor $H-F(\rho)$ measures the kinetic part of the energy. Thus the
potential does not replace the Wasserstein metric; it changes the geometry by
making some regions of probability space more expensive than others.

\begin{remark}
This is the first point where the bridge stops looking like ordinary optimal
transport. The geometry is no longer only the bare Wasserstein geometry of
mass movement. It is Wasserstein geometry shaped by an energy landscape.
\end{remark}

\section{Why Conservative Bridges Are Not Enough}

The Hamiltonian formulation above is conservative: the total Hamiltonian is
constant along the idealized flow. This is appropriate for systems whose
energy is preserved, but many systems of interest in generative modeling and
reinforcement learning are not like that.

Examples include:

\begin{itemize}
    \item a policy distribution changing as new value information is injected,
    \item a generative model being guided toward a condition,
    \item a biological population differentiating over time,
    \item an offline policy being updated toward an online-improved policy.
\end{itemize}

In these cases, the path may gain or lose effective energy. The target itself
may drift. The geometry should therefore allow the bridge to be
non-conservative.

\section{The Non-Conservative Lift}

The non-conservative bridge introduces one extra scalar variable $z_t$. This
variable stores the accumulated action along the path. Instead of minimizing a
fixed action directly, one prescribes its evolution:

\[
\partial_t z_t
=
\int_{\R^d}
\left(
\frac12\|v_t(x)\|^2+U(x)
\right)
\rho_t(x)\,dx
-\gamma z_t.
\]

The parameter $\gamma\in\R$ controls damping or amplification. The bridge is
still constrained by the Fokker--Planck equation

\[
\partial_t\rho_t+\nabla\cdot(\rho_t v_t)
=
\varepsilon\Delta\rho_t.
\]

\begin{definition}[Non-conservative generalized Schr\"odinger bridge]
A non-conservative generalized Schr\"odinger bridge is a controlled density
path $(\rho_t,v_t,z_t)$ satisfying

\[
\partial_t\rho_t+\nabla\cdot(\rho_t v_t)
=
\varepsilon\Delta\rho_t
\]

and

\[
\partial_t z_t
=
\int_{\R^d}
\left(
\frac12\|v_t(x)\|^2+U(x)
\right)
\rho_t(x)\,dx
-\gamma z_t,
\]

with prescribed endpoint marginals, and possibly prescribed intermediate
marginals.
\end{definition}

The point of $z_t$ is not that probability mass is no longer conserved. The
mass is still conserved by the Fokker--Planck equation. What is no longer
conserved is the energy/action budget governing the bridge.

\section{Contact Hamiltonian Dynamics}

Contact geometry is the odd-dimensional analogue of symplectic geometry. The
usual Hamiltonian phase space is enlarged by one scalar coordinate. In our
case, the enlarged phase space is formally

\[
T^*\cP^+(\R^d)\times\R,
\]

with coordinates

\[
(\rho,S,z).
\]

The contact Hamiltonian has the schematic form

\[
H^c(\rho,S,z)
=
K(\rho,S)+F(\rho)+B(\rho)+\gamma z.
\]

Here $K$ is kinetic energy, $F$ contains the potential and diffusion-induced
terms, $B$ is an entropy contribution, and $\gamma z$ is the explicitly
non-conservative part.

The contact Hamiltonian equations have the form

\[
\partial_t\rho_t
=
-\nabla\cdot(\rho_t\nabla S_t),
\]

\[
\partial_t S_t
=
-\frac12\|\nabla S_t\|^2
+\frac12\varepsilon^2
\frac{\delta I}{\delta\rho}(\rho_t)
+U
-\gamma S_t
-\varepsilon\gamma\log\rho_t,
\]

and

\[
\partial_t z_t
=
\int
\left(
\frac12\|\nabla S_t(x)\|^2+U(x)
\right)
\rho_t(x)\,dx
+\frac12\varepsilon^2 I(\rho_t)
-\int
\varepsilon\gamma(\log\rho_t(x)-1)\rho_t(x)\,dx
-\gamma z_t.
\]

\begin{remark}
The terms involving $\gamma$ are the signature of non-conservative dynamics.
When $\gamma=0$, the extra contact effects disappear and the system reduces
to the conservative Hamiltonian bridge picture.
\end{remark}

\section{Contact Wasserstein Geodesics}

The contact Hamiltonian flow lives on the extended phase space
$T^*\cP^+(\R^d)\times\R$. Its projection can be interpreted as a geodesic on an
augmented probability manifold

\[
\cP^+(\R^d)\times\R.
\]

The associated metric is again a Jacobi-type rescaling of the Wasserstein
metric:

\[
\widetilde g_J
=
\bigl(H^c-F(\rho)-B(\rho)\bigr)g_{W_2}.
\]

This metric is still Wasserstein in its tangent directions, but the cost of
moving through probability space is modulated by the evolving
non-conservative energy.

\begin{definition}[Contact Wasserstein geodesic]
A contact Wasserstein geodesic is a probability path, together with an
accumulated action coordinate $z_t$, that solves the non-conservative bridge
problem and is equivalently described as a geodesic for the contact
Jacobi-type metric on $\cP^+(\R^d)\times\R$.
\end{definition}

\section{Finite-Dimensional Parameterization}

The full space $\cP^+(\R^d)$ is infinite-dimensional. To compute bridges in
practice, one restricts attention to a finite-dimensional family of
pushforward distributions. Let $\lambda$ be a reference distribution and let

\[
T_{\theta_K}\circ\cdots\circ T_{\theta_1}\circ T_{\theta_0}
\]

be a composition of maps. This induces a discrete sequence

\[
\rho^{\theta}_{t_{k+1}}
=
(T_{\theta_{k+1}})_\#
\rho^{\theta}_{t_k}.
\]

This is exactly the geometric role played by a residual network: each block is
a small transport step, and the whole network is a discretized path in
probability space.

The training loss has three conceptual parts:

\[
\ell
=
\underbrace{
d_{W_2}^2(\rho^\theta_{t_K},\rho_b)
}_{\text{terminal marginal}}
+
\underbrace{
\sum_{m=1}^M
d_{W_2}^2(\rho^\theta_{t_{k_m}},\rho_m)
}_{\text{intermediate marginals}}
+
\underbrace{
\sum_{k=1}^K
\Phi_{t_k}
d_{W_2}^2(\rho^\theta_{t_k},\rho^\theta_{t_{k-1}})
}_{\text{contact geodesic energy}}.
\]

The scalar $\Phi_{t_k}$ is the discrete version of the contact Jacobi factor.
It tells the model how expensive it is to move through that portion of the
probability path.

\begin{remark}
This is the computational bridge between the geometry and generative modeling.
The abstract object is a geodesic in an infinite-dimensional Wasserstein-type
space. The implemented object is a sequence of parameterized pushforwards
trained to match marginals while minimizing a geometry-weighted transport
energy.
\end{remark}

\section{Guidance as Metric Modulation}

Suppose we want the bridge to satisfy a condition

\[
y=f(x_{t_s})
\]

at some time $t_s$. A natural way to impose this condition is to add the
penalty

\[
\|y-f(x_{t_s})\|^2
\]

to the action. Geometrically, this modifies the Jacobi factor:

\[
\widetilde g'_J
=
\left(
\Phi_t+\|y-f(x_{t_s})\|^2
\right)g_{W_2}.
\]

Regions of probability space that violate the condition become more expensive.
The bridge is therefore steered without abandoning the underlying transport
geometry.

\section{Interpretation for Generative Policies}

For generative policies, a policy is a pushforward

\[
\pi_\theta(\cdot|s)
=
f_\theta(s,\cdot)_\#\rho_0.
\]

Training changes the induced distribution over actions. The conservative
Wasserstein-gradient-flow picture says that policy improvement is a
mass-preserving transport process driven by a variational objective.

The contact-Wasserstein picture suggests a richer possibility:

\[
\pi_\beta
\longrightarrow
\pi_Q
\]

may be better modeled as a non-conservative bridge whose effective energy can
change during learning.

In this interpretation:

\begin{itemize}
    \item the behavior policy is the initial marginal,
    \item the improved soft policy is the terminal marginal,
    \item value information acts like a potential or guidance term,
    \item entropy controls exploration,
    \item trust-region effects can be encoded as transport costs,
    \item offline-to-online updates can inject or dissipate energy.
\end{itemize}

This is the conceptual reason for studying contact Wasserstein geometry in
the context of drifting policies. It gives a language for distributional
learning dynamics that are still geometric, but no longer forced to behave as
closed conservative systems.
